\subsection{Ramification and rational-codegree trace bounds}
\label{subsec:ramification-trace}
This criterion applies to arbitrary complex fusion categorifications.  It
does not assume a pivotal structure and does not enumerate selected
dimension characters.  Let $R$ be commutative and self-dual of rank six,
and write $T_0=\sum_{i,j,k}(N_{ij}^k)^2$.

\begin{proposition}[Ramified dimension field and Casimir trace]
\label{prop:ramification-trace}
Suppose the FP global dimension is
$D=a+b\sqrt q$, where $q>1$ is squarefree,
$q\equiv2,3\pmod4$, and $a,b\in\ZZ$, $b>0$.
Suppose also that
\[
 R/2R\cong\FF_2[z]/(z^2)\times\prod_j\FF_{2^{e_j}}.
\]
Let $\mathcal L\subseteq\ZZ^2$ be the lattice of pairs $(S,B)$ for which
\[
 \frac{2qaS-2qbB}{a^2-qb^2},\qquad
 \frac{2aB-2qbS}{a^2-qb^2},\qquad
 \frac{2qS}{c},\qquad\frac{2B}{c}
 \quad\hbox{are integers}
\]
for every rational original formal codegree $c$.
Let $v$ be the positive generator of the projection of $\mathcal L$
onto its first coordinate.  If
\[
 4qv+5>T_0,
\]
then $R$ has no complex fusion categorification.
\end{proposition}
\begin{proof}
Sphericalize a hypothetical categorification and quotient its Grothendieck
ring by the central sign object plus the unit.  The result is a rank-six
integral signed algebra $A$ with multiplication coefficients $t_{ij}^k$
satisfying
\[
 t_{ij}^0=\delta_{ij},\qquad
 |t_{ij}^k|\le N_{ij}^k,\qquad
 t_{ij}^k\equiv N_{ij}^k\pmod2.
\]
The two lifts of each original self-dual simple have opposite nonzero
spherical dimensions, so duality cannot interchange them.  Consequently
the regular matrices $M_i$ of $A$ are real symmetric and commute.
The dimension character has values $d_i$ in a totally real field factor
$K$ of $A\otimes\QQ$, and $\Delta=\sum_i d_i^2$ is the global dimension
of the original category.

For every original formal codegree $f$, its pulled-back character on
sphericalization has formal codegree $2f$.
The central-induction theorem
\cite[Theorem~2.13]{Ostrik3} therefore gives a central simple of dimension
$\Delta/f$.  Forgetting it expresses this dimension as an integer
combination of the $\pm d_i$.  Thus $\Delta/f\in\mathcal O_K$, and in
particular $D\in K$ and $\QQ(\sqrt q)\subseteq K$.
This also recalls the dimension-field argument of
Lemma~\ref{lem:dimension-field-containment}.

Since $A/2A=R/2R$, lift the displayed primitive idempotents to
$A\otimes\ZZ_2$.  A finite free factor with finite-field reduction is
finite etale over $\ZZ_2$, so its generic fiber is unramified.  The only
remaining factor has rank two.  But $\QQ_2(\sqrt q)$ is ramified for
$q\equiv2,3\pmod4$, and every component of $K\otimes\QQ_2$ contains
this field.  No such component can be supported on an unramified factor.
Hence $[K:\QQ]\le2$, and therefore $K=\QQ(\sqrt q)$.

Now $\mathcal O_K=\ZZ[\sqrt q]$.  Write
$d_i=\alpha_i+\beta_i\sqrt q$ for $1\le i\le5$, and set
$S=\sum_i\beta_i^2>0$ and $B=\sum_i\alpha_i\beta_i$.
The two conjugate dimension-character vectors are orthogonal, because
the $M_i$ are commuting real symmetric matrices.  It follows that
\[
 1+\sum_i\alpha_i^2=qS,\qquad
 \Delta=2qS+2B\sqrt q,\qquad
 \Delta+\Delta'=4qS.
\]
Integrality of $\Delta/D$ gives the first two conditions defining
$\mathcal L$.  Integrality of $\Delta/c$ gives the other two for every
rational codegree $c$.  Therefore $S\ge v$.

The signed Casimir $\sum_iM_i^2$ has the six signed formal codegrees as
its eigenvalues.  The two dimension characters contribute $4qS$.
Each of the other four codegrees is strictly greater than one: equality
would force every nonunit character value to vanish, contrary to the
unit coefficient in $x_i^2$.  Since the trace is integral,
\[
 4qv+5\le\Tr\Bigl(\sum_iM_i^2\Bigr)
 =\sum_{i,j,k}(t_{ij}^k)^2\le T_0,
\]
contradicting the asserted strict inequality.
\end{proof}

\begin{corollary}
The rings $\RR{7}{3812}$, $\RR{8}{6234}$, and $\RR{8}{6509}$ have no
complex fusion categorifications.
\end{corollary}
\begin{proof}
For all three rings,
\[
 q=226,\qquad D=2260+150\sqrt{226},\qquad
 D^2-4520D+22600=0.
\]
The following table lists their exact local factors and trace data.
Here $B_2=\FF_2[z]/(z^2)$ and $D'$ is the conjugate of $D$.
\begin{center}
\begin{tabular}{@{}llllr@{}}
\toprule
Ring & Algebra modulo two & Formal codegrees & $v$ & $T_0$\\
\midrule
$\RR{7}{3812}$ & $B_2\times\FF_4\times\FF_2^2$
 & $D,D',3,15,5,5$ & 15 & 4548\\
$\RR{8}{6234}$ & $B_2\times\FF_4\times\FF_2^2$
 & $D,D',3,15,5,5$ & 15 & 4548\\
$\RR{8}{6509}$ & $B_2\times\FF_2^4$
 & $D,D',3,3,15,15$ & 15 & 4556\\
\bottomrule
\end{tabular}
\end{center}
In particular $\Delta/15\in\mathcal O_K$ and
$\Delta=452S+2B\sqrt{226}$ imply $15\mid S$, because
$\gcd(452,15)=1$.  Thus the signed Casimir trace is at least
$904\cdot15+5=13565$, exceeding each displayed $T_0$.
\end{proof}

The finite local checks exhaust all $64$ elements modulo two, find every
primitive idempotent, and verify orthogonality and completeness of the
factor decomposition.  Every nonzero element of a claimed field factor
is checked to have an inverse; the remaining rank-two factor has exactly
one nonzero square-zero element.  The lattice step is recovered from exact
congruences, not floating-point bounds.  The three certificates require
neither a signed-table search nor a pseudo-unitary exclusion.

The tensor-based implementation is
\filename{software/new_ramified_signed.py}.  It was applied to the frozen
second-round frontier of $40$ rings, including all $39$ with multiplicity
four through eight.  Of those $39$, the hypotheses fail for $34$, one
($\RR{7}{3582}$) proceeds to a bounded signed reduction, and four are
excluded directly: the three above and an independent field-containment
exclusion of $\RR{8}{6940}$.  The latter is also covered by the
biquadratic criterion, so it is not counted twice.
The complete per-input report is
\filename{research/round2_higher/ramification_application.json}.
All $54$ positive models are preserved.  Ordinary and optimized Python
checks verify the three certificates after nontrivial basis permutations
and reject six certificate corruptions and a corrupted tensor.
